\section{Classical inputs}\label{app:classical}

This appendix collects the routine parts of Sections~\ref{sec:setting}
and~\ref{sec:bulk}. They are the proofs of Lemmas~\ref{lem:rv}
and~\ref{lem:profile}, four classical facts on sums of runs, and the bounds on
the kernels of one sign that enter the proof of Theorem~\ref{thm:bulk}.

\subsection{Regular variation and the profile}\label{app:rv}

\begin{proof}[Proof of Lemma~\ref{lem:rv}]
Parts (i) to (iii) are the cited theorems, and $T_\sgn(x/c)\le C_cT_\sgn(x)$ follows
from (i) for large $x$ and from the positivity of the nonincreasing $T_\sgn$ for the
others. For (iv), $nT_\sgn(n)$ has index $1-\alpha_\sgn>-1$ and $T_\sgn$ has index
$-\alpha_\sgn<-1$, so the cited propositions give the asymptotics, and positivity gives
the bounds for all $w\ge1$. For (v), each $a_\sgn$ is increasing, so $B_N$ is
nondecreasing. If $\alpha_\plus\ne\alpha_\minus$, the ratio $b_{\plus,N}/b_{\minus,N}$
is regularly varying with a nonzero index, so $B_N=b_{\sgn,N}$ for the sign with the
smaller index and all large $N$. If the indices are equal, $B_{cN}/B_N$ lies between
the two ratios $b_{\sgn,cN}/b_{\sgn,N}$, which tend to $c^{1/\alpha_*}$. Then
$B_NT_\sgn(N)$ has index $1/\alpha_*-\alpha_\sgn<0$, and a regularly varying function
with a negative index tends to zero~\cite[Proposition~1.5.1]{bgt1989}. The definition of
$a_\sgn$ with $t=N/s$ gives $NT_\sgn(b_{\sgn,N})\to s$, and $B_N\log N=o(N)$ because
$1/\alpha_*<1$. For (vi), take $\delta=(\alpha_\sgn-1)/2\le\frac12$ in (ii). For large
$N$ we have $b_{\sgn,N}\ge x_\delta$ and $s/2\le NT_\sgn(b_{\sgn,N})\le2s$, and
$(L/b_{\sgn,N})^{-\alpha_\sgn-\delta}\ge(L/b_{\sgn,N})^{-\alpha_\sgn-1}$ for $L\ge b_{\sgn,N}$.
\end{proof}

\begin{proof}[Proof of Lemma~\ref{lem:profile}]
\emph{Bounds and continuity.} Under (RV) the means are finite, so $0<p<1$. For
$1<\alpha<2$ and the principal branch, $(\mp it)^\alpha=|t|^\alpha e^{\mp i\pi\alpha\,\mathrm{sgn}(t)/2}$, so
the real parts of $(-it)^\alpha$ and $(it)^\alpha$ equal $|t|^\alpha\cos\frac{\pi\alpha}2<0$, and
$(-ict)^\alpha=c^\alpha(-it)^\alpha$ for $c\ge0$. Hence $W_\theta$ has the
characteristic function
\[
 \widehat h_\theta(t)=\exp\Bigl(K_{\alpha_\plus}\bigl((1-p)\theta_\plus\bigr)^{\alpha_\plus}(-it)^{\alpha_\plus}
 +K_{\alpha_\minus}\bigl(p\,\theta_\minus\bigr)^{\alpha_\minus}(it)^{\alpha_\minus}\Bigr),
\]
which is continuous in $\theta$ for each $t$. With
$c_\sgn=K_{\alpha_\sgn}\lvert\cos\frac{\pi\alpha_\sgn}2\rvert$ and
$c=\min(c_\plus(1-p)^{\alpha_\plus},c_\minus p^{\alpha_\minus})>0$, one coordinate of
$\theta\in\Theta$ equals $1$, so $|\widehat h_\theta(t)|\le\Phi(t)$, where $\Phi(t)=1$ for
$|t|<1$ and $\Phi(t)=e^{-c|t|^{\alpha_*}}$ for $|t|\ge1$. Both $\Phi$ and $|t|\Phi(t)$
are integrable. Fourier inversion gives the continuous density
$h_\theta(x)=\frac1{2\pi}\int e^{-itx}\widehat h_\theta(t)\,dt$, with
$\lVert h_\theta\rVert_\infty\le C_1=\frac1{2\pi}\int\Phi$ and Lipschitz constant
$C_2=\frac1{2\pi}\int|t|\Phi(t)\,dt$, and $h_\theta$ vanishes at infinity by the
Riemann--Lebesgue lemma. This gives (a). Also
$\lVert h_\theta-h_{\theta'}\rVert_\infty\le\frac1{2\pi}\int|\widehat h_\theta-\widehat h_{\theta'}|$,
and dominated convergence with the bound $2\Phi$ gives the continuity in (c). If
$N'/N\to1$, then $a_\sgn(N'/s)/a_\sgn(N/s)\to1$ by Lemma~\ref{lem:rv}(i), so
$\theta_{N'}-\theta_N\to0$, and the uniform continuity of $\theta\mapsto h_\theta$ on the
compact $\Theta$ gives the last claim of (c).

\emph{Tails.} The set $\{h_\theta:\theta\in\Theta\}$ is compact. Given $\varepsilon>0$
take a finite $\varepsilon$-net $h_{\theta_1},\dots,h_{\theta_q}$ of it and $R$ with
$h_{\theta_i}<\varepsilon$ on $|x|\ge R$ for all $i$. Then $h_N<2\varepsilon$ on $|x|\ge R$
for every $N$, so $\varepsilon_h(R)\to0$. If $\varepsilon'_h(R)\not\to0$, there are
$\varepsilon>0$, $R_i\to\infty$ and $N_i$ with $\int_{|x|>R_i}h_{N_i}\ge\varepsilon$, and by
compactness we may assume $\theta_{N_i}\to\theta$. Then $h_{N_i}\to h_\theta$ uniformly,
so in $L^1$ by Scheff\'e's lemma, and
$\int_{|x|>R_i}h_{N_i}\le\lVert h_{N_i}-h_\theta\rVert_1+\int_{|x|>R_i}h_\theta\to0$, a
contradiction. This proves (b).

\emph{Positivity.} Fix a sign, write $Z=Z_\sgn$, $g=g_\sgn$ and $\alpha=\alpha_\sgn$, and let
$Z_1,Z_2,\dots$ be independent copies of $Z$. The scaling of the characteristic
function shows that $Z_1+\dots+Z_n$ has the law of $n^{1/\alpha}Z$. If $Z\ge0$ almost
surely, then $\delta=\E\min(Z,1)>0$, and by the strong law of large numbers
$n^{-1/\alpha}(Z_1+\dots+Z_n)\ge n^{1-1/\alpha}\cdot\frac1n\sum_{i\le n}\min(Z_i,1)\to\infty$
almost surely, which contradicts the tightness of a sequence with the law of $Z$.
The same holds for $-Z$. So the open set $U=\{g>0\}$ meets both half lines. The
density $g^{*n}$ of $Z_1+\dots+Z_n$ is continuous and positive at $y_1+\dots+y_n$
whenever $y_1,\dots,y_n\in U$, and it equals $n^{-1/\alpha}g(n^{-1/\alpha}\,\cdot)$. Hence
$y_1,\dots,y_n\in U$ implies $n^{-1/\alpha}(y_1+\dots+y_n)\in U$. Take
$I=(a-\varepsilon,a+\varepsilon)\subset U$ with $0<\varepsilon<a$ and put $c_n=n^{1-1/\alpha}$.
Then $c_nI\subset U$ for every $n$. Since $c_n\to\infty$ and $c_{n+1}/c_n\to1$, the
intervals $c_nI$ with $n\ge n_0$ overlap in turn, so $U$ contains a half line
$(x_0,\infty)$. In the same way $U$ contains $(-\infty,-x_0)$. Every real number is the
sum of a point of each half line, so the case $n=2$ gives $U=\R$. Now let $\theta\in\Theta$
and $\theta_\sgn=1$. Then $W_\theta=aZ_\sgn+Y$ with $a\in\{1-p,-p\}$ and $Y$ independent
of $Z_\sgn$, so $x\mapsto\E[|a|^{-1}g_\sgn((x-Y)/a)]$ is a continuous positive density of
$W_\theta$, and it equals $h_\theta$. The maps $\theta\mapsto\sup_xh_\theta(x)$ and
$\theta\mapsto\min_{|x|\le K}h_\theta(x)$ are continuous and positive on the compact
$\Theta$, so their minima are positive. This proves (d).

\emph{The formula.} Let $\tau=\plus$. Given $Z_\minus=y$, the variable
$(1-p)Z_\plus-p\theta_{\minus,N}Z_\minus$ has the density
$(1-p)^{-1}g_\plus((x+p\theta_{\minus,N}y)/(1-p))$ at $x$. Since $1-p=\mu_\minus/s$ and
$p/(1-p)=\mu_\plus/\mu_\minus$, integrating against $g_\minus(y)\,dy$
gives~\eqref{eq:bulk-profile}. For $\tau=\minus$ we condition on $Z_\plus$.
\end{proof}

\subsection{Sums of runs}\label{sec:bulk-inputs}

We use four facts about the sums $S^\sgn_m$, stated with the normalization of
Section~\ref{sec:scales}. We write $\bar T_\sgn(x)=\Prob(X_\sgn>x)$, which is
regularly varying with $\bar T_\sgn(x)\sim T_\sgn(x)$ under (RV).

\begin{lemma}[stable and local limits]\label{lem:stable}
Assume (RV) and (SPAN) for the sign $\sgn$, and write $\alpha=\alpha_\sgn$,
$\mu=\mu_\sgn$, $a=a_\sgn$, $g=g_\sgn$, $T=T_\sgn$ and $S_m=S^\sgn_m$. Parts
(i), (iii) and (iv) need only (RV).
\begin{enumerate}[(i)]
\item $Z_m=(S_m-m\mu)/a(m)$ converges in law to $Z_\sgn$, whose density $g$ is
bounded, uniformly continuous, and tends to zero at infinity. Moreover
$m\bigl(\E e^{-(X_\sgn-\mu)/a(m)}-1\bigr)\to K_\alpha$.
\item (Gnedenko's local limit theorem)
$\sup_{k\in\Z}\bigl|a(m)\Prob(S_m=k)-g\bigl((k-m\mu)/a(m)\bigr)\bigr|\to0$ as $m\to\infty$,
and $\bar u_\sgn=\sup_{m\ge1}\sup_ka(m)\Prob(S_m=k)<\infty$.
\item (local large deviations) There is a constant $C_{\mathrm{LD}}=C_{\mathrm{LD}}(\sgn)$ such that
$\Prob(S_m=j)\le C_{\mathrm{LD}}\,m\,T(j-m\mu)/a(m)$
for all $m\ge1$ and all integers $j$ with $j-m\mu\ge1$.
\item (left deviations) The constant $C_-=C_-(\sgn)=\sup_{m\ge1}m\log\E e^{-(X_\sgn-\mu)/a(m)}$
is finite, and
$\Prob(S_m\le m\mu-y)\le e^{C_-}\,e^{-y/a(m)}$ for all $m\ge1$ and $y\ge0$.
\end{enumerate}
\end{lemma}

\begin{proof}
(i) The domain of attraction is classical~\cite[Section~XVII.5]{feller1971}, and we
check the normalization, since the constant $K_\alpha$ is ours. Let $X=X_\sgn$ and, for
$\operatorname{Re}\xi\ge0$, $\Lambda(\xi)=\E[e^{-\xi X}-1+\xi X]$. Since
$e^{-\xi x}-1+\xi x=\int_0^x\xi(1-e^{-\xi y})\,dy$ for $x\ge0$, Fubini's theorem gives,
with $\xi=te^{i\vartheta}$, $t>0$ and $|\vartheta|\le\pi/2$,
\[
 \Lambda(\xi)=\int_0^\infty \xi\bigl(1-e^{-\xi y}\bigr)\bar T(y)\,dy
 =\int_0^\infty e^{i\vartheta}\bigl(1-e^{-e^{i\vartheta}v}\bigr)\bar T(v/t)\,dv .
\]
The integrand is at most $\min(2,v)\bar T(v/t)$ in modulus. As $t\to0$,
$\bar T(v/t)/\bar T(1/t)\to v^{-\alpha}$, and Lemma~\ref{lem:rv}(ii) dominates the ratio
by $2\max(v^{-\alpha-\delta},v^{-\alpha+\delta})$ for $v\ge tx_\delta$, with
$1<\alpha\pm\delta<2$. The part $v<tx_\delta$ contributes at most
$t^2x_\delta^2/\bar T(1/t)\to0$ to $\Lambda(\xi)/\bar T(1/t)$, because $\alpha<2$. Hence,
uniformly in $\vartheta$,
\[
 \Lambda(te^{i\vartheta})=\bar T(1/t)\bigl(F(e^{i\vartheta})+o(1)\bigr),\qquad
 F(\xi)=\int_0^\infty \xi\bigl(1-e^{-\xi v}\bigr)v^{-\alpha}\,dv .
\]
The function $F$ is continuous on $\operatorname{Re}\xi\ge0$, analytic inside, and
$F(t)=t^\alpha F(1)$ for $t>0$. An integration by parts gives
$F(1)=\Gamma(2-\alpha)/(\alpha-1)=K_\alpha$, so $F(\xi)=K_\alpha\xi^\alpha$ with the principal
branch. Next $\E e^{-\xi(X-\mu)}-1=\Lambda(\xi)+O(|\xi|^2)$ as $\xi\to0$. With
$\xi=-i\theta/a(m)$, $m\bar T(a(m)/|\theta|)\to|\theta|^\alpha$ and $m/a(m)^2\to0$ give
$m(\E e^{i\theta(X-\mu)/a(m)}-1)\to K_\alpha(-i\theta)^\alpha$, and L\'evy's continuity
theorem gives $Z_m\to Z_\sgn$. The characteristic function of $Z_\sgn$ has modulus
$\exp(K_\alpha\cos(\pi\alpha/2)|\theta|^\alpha)$, which is integrable, so Fourier inversion
and the Riemann--Lebesgue lemma give the properties of $g$. With $\xi=1/a(m)$ the same
expansion gives the last claim.

(ii) The support of $X_\sgn-\mu$ lies in $\Z-\mu$, and by (SPAN) its maximal lattice
span is $1$. Given (i), Gnedenko's local limit
theorem~\cite[Theorem~4.2.1]{ibragimovlinnik1971} (see
also~\cite[Theorem~8.4.1]{bgt1989}) gives the uniform limit. The norming sequence
there may be any $a_m\sim a(m)$, since $g$ is bounded, uniformly continuous and tends
to zero at infinity. For $\bar u_\sgn$, the limit gives an $m_0$ with
$a(m)\Prob(S_m=k)\le\|g\|_\infty+1$ for $m\ge m_0$, and $a(m)\Prob(S_m=k)\le a(m_0)$ for
$m<m_0$.

(iii) Caravenna and Doney~\cite[Theorem~1.1]{caravennadoney2019} assume right and
left tails $p/A(x)$ and $q/A(x)$ (or $o(1/A(x))$) with $A$ regularly varying of index
in $(0,1)\cup(1,2)$, $p>0$, and mean zero when the index exceeds one. Their bound
(1.6) then holds with no lattice assumption, for all $m\ge1$, $x\ge0$ and bounded
intervals $\mathcal I$, and reads
$\Prob(S_m-m\mu\in x+\mathcal I)\le C_0'(\mathcal I)\,m/(a_mA(x))$ with $a_m=A^{-1}(m)$.
Under (RV) the variable $X_\sgn-\mu$ satisfies these hypotheses with $p=1$, $q=0$ and
$A(x)\sim1/T(x)$. We take $\mathcal I=(-1,0]$ and $x=j-m\mu$, so that
the event contains $\{S_m=j\}$. The asymptotic inverses $a_m$ and $a(m)$ satisfy
$a_m/a(m)\to1$~\cite[Theorem~1.5.12]{bgt1989}, so $a(m)\le c\,a_m$ for all $m$, and
$1/A(x)\le c\,T(x)$ for $x\ge1$. This gives (iii).

(iv) By (i), $m(\E e^{-(X_\sgn-\mu)/a(m)}-1)$ is bounded, and $\log(1+z)\le z$ makes
$C_-$ finite. Markov's inequality for $e^{-(S_m-m\mu)/a(m)}$ and independence give
$\Prob(S_m-m\mu\le-y)\le e^{-y/a(m)}(\E e^{-(X_\sgn-\mu)/a(m)})^m\le e^{C_--y/a(m)}$.
This is the exponential bound behind the Fuk--Nagaev inequalities~\cite{nagaev1979}.
\end{proof}

\subsection{Kernels of one sign}\label{sec:bulk-kernels}

For each sign $\sgn$ the fresh identities use three kernels,
\[
 \mathsf k^{\sgn,1}_m=u^\sgn_m,\qquad \mathsf k^{\sgn,2}_m=q^\sgn_m,\qquad
 \mathsf k^{\sgn,3}_m=q^\sgn_{m+1}\qquad(m\ge1),
\]
with weights $c_{\mathsf k^{\sgn,1}}=1$ and $c_{\mathsf k^{\sgn,2}}=c_{\mathsf k^{\sgn,3}}=\mu_\sgn$. A
kernel of sign $\sgn$ is one of these three sequences, and we write $\mathsf k_m$ and
$c_{\mathsf k}$ when the choice does not matter. We use four facts.
\begin{enumerate}[(K1)]
\item $\sum_{m\ge1}u^\sgn_m(j)\le1$, $\sum_{m\ge1}q^\sgn_m(j)=1$ and
$\sum_{m\ge1}q^\sgn_{m+1}(j)\le1$ for $j\ge1$, by~\eqref{eq:kernel-sums}.
\item $\mathsf k_m(j)=0$ for $m>j$, because $S^\sgn_m\ge m$.
\item Let $J_\sgn(j)=\min\{m\ge1:S^\sgn_m\ge j\}$. Then $\Prob(J_\sgn(j)=m)=q^\sgn_m(j)$
by~\eqref{eq:q-cover}, and $u^\sgn_m(j)\le q^\sgn_m(j)$, since $\{S^\sgn_m=j\}\subset\{J_\sgn(j)=m\}$.
\item $\sum_jq^\sgn_m(j)=\mu_\sgn$ and $\sum_ju^\sgn_m(j)=1$.
\end{enumerate}

\begin{lemma}[kernel local limits]\label{lem:kernels}
Assume (RV) and (SPAN). For each sign $\sgn$ and each kernel $\mathsf k$ of sign $\sgn$,
\[
 \delta_{\mathsf k}(m)=\sup_{j\in\Z}\bigl|a_\sgn(m)\mathsf k_m(j)-c_{\mathsf k}\,g_\sgn\bigl((j-m\mu_\sgn)/a_\sgn(m)\bigr)\bigr|\to0
 \qquad(m\to\infty),
\]
and $\mathsf k_m(j)\le C_3/a_\sgn(m)$ for all $m\ge1$ and $j$, where
$C_3=C_3(\sgn)=\max\bigl(a_\sgn(1),\,\mu_\sgn r_\sgn\bar u_\sgn\bigr)$ and
$r_\sgn=\sup_{m\ge2}a_\sgn(m)/a_\sgn(m-1)<\infty$.
\end{lemma}

\begin{proof}
For $u^\sgn_m$ this is Lemma~\ref{lem:stable}(ii). The kernel $q^\sgn_m$ is a sum of
shifts of $u^\sgn_{m-1}$ with weights $T_\sgn(v)$ of total mass $\mu_\sgn$, and
$q^\sgn_{m+1}$ is the same sum built on $u^\sgn_m$. This gives the bounds, with
$q^\sgn_1=T_\sgn\le1$ at $m=1$. For the limit of $q^\sgn_m$ write $a=a_\sgn$, $g=g_\sgn$,
$z=(j-m\mu)/a(m)$, and fix $v_0\ge1$ and $Z>0$. The shifts $v>v_0$ contribute at most
$(\bar u_\sgn+\|g\|_\infty)\sum_{v>v_0}T(v)$ to $|a(m-1)q_m(j)-\mu g(z)|$. For $v\le v_0$
the point $(j-v-(m-1)\mu)/a(m-1)=z\,a(m)/a(m-1)+(\mu-v)/a(m-1)$ lies within
$\epsilon_m=Z|a(m)/a(m-1)-1|+(v_0+\mu)/a(m-1)$ of $z$ when $|z|\le Z$, and has modulus at
least $Z/2$ for large $m$ when $|z|>Z$. So for large $m$
\[
 \bigl|a(m-1)q_m(j)-\mu g(z)\bigr|\le\mu\bigl(\delta_{u}(m-1)+\omega_g(\epsilon_m)+\bar g(Z/2)\bigr)
 +(\bar u_\sgn+\|g\|_\infty)\sum_{v>v_0}T(v),
\]
where $\omega_g$ is the modulus of continuity of $g$ and $\bar g(r)=\sup_{|y|\ge r}g(y)$.
Let $m\to\infty$, then $v_0,Z\to\infty$, and use $a(m)/a(m-1)\to1$. The kernel
$q^\sgn_{m+1}$ is the same with the shift $-v/a(m)$ in place of $(\mu-v)/a(m-1)$.
\end{proof}

\begin{lemma}[small counts]\label{lem:small}
Assume (RV). For each sign $\sgn$ there is $C_4=C_4(\sgn)$ such that every kernel
of sign $\sgn$ satisfies
\[
 \mathsf k_m(j)\le C_4\bigl(1+m/a_\sgn(m)\bigr)T_\sgn(j)\qquad\text{whenever }8m\mu_\sgn\le j .
\]
Consequently, for $0<\beta\le1$ there is $C_5=C_5(\sgn,\beta)$ with
\begin{equation}\label{eq:bulk-small}
 \mathsf k_m(j)\le C_5\,e_{\sgn,N},\qquad e_{\sgn,N}=\bigl(1+N/a_\sgn(N)\bigr)T_\sgn(N),
\end{equation}
for all $N\ge1$, all $\beta N\le j\le N$ and all $m$ with $8m\mu_\sgn\le\beta N$.
Moreover $B_Ne_{\sgn,N}\to0$.
\end{lemma}

\begin{proof}
Write $T=T_\sgn$, $a=a_\sgn$, $\mu=\mu_\sgn$, and let $C_T$ be the constant of
Lemma~\ref{lem:rv}(ii) with $T(x/2)\le C_TT(x)$. If $4m\mu\le i$, then
$i-m\mu\ge i/2\ge1$, and Lemma~\ref{lem:stable}(iii) gives
\begin{equation}\label{eq:bulk-u-small}
 u_m(i)\le C_{\mathrm{LD}}\,m\,T(i/2)/a(m)\le C_{\mathrm{LD}}C_T\,m\,T(i)/a(m).
\end{equation}
This is the claim for $u_m$. Let $8m\mu\le j$. For $q_1=T$ there is nothing to prove.
For $q_m$ with $m\ge2$ we split $\sum_vT(v)u_{m-1}(j-v)$ at $v=j/2$. If $v\le j/2$, then
$j-v\ge j/2\ge4(m-1)\mu$, so~\eqref{eq:bulk-u-small} bounds $u_{m-1}(j-v)$ by
$C_{\mathrm{LD}}C_T^2r_\sgn\,m\,T(j)/a(m)$, and the weights $T(v)$ add up to at most
$\mu$. If $v>j/2$, then $T(v)\le C_TT(j)$, and the values $u_{m-1}(j-v)$ add up to at
most $1$ by (K4). So $q_m(j)\le(\mu C_{\mathrm{LD}}C_T^2r_\sgn\,m/a(m)+C_T)T(j)$. The
kernel $q_{m+1}$ is the same with $u_m$ in place of $u_{m-1}$. This proves the first
claim with $C_4=C_T+\mu C_{\mathrm{LD}}C_T^2r_\sgn$.

For~\eqref{eq:bulk-small}, $m/a(m)$ has the positive index $1-1/\alpha_\sgn$, so
Lemma~\ref{lem:rv}(iii) gives $m/a(m)\le C_PN/a(N)$ for $m\le N$, and
$T(j)\le T(\beta N)\le C_{T,\beta}T(N)$ by Lemma~\ref{lem:rv}(ii). Finally
$B_NNT_\sgn(N)/a_\sgn(N)$ has the index
\[
 \frac1{\alpha_*}-\frac1{\alpha_\sgn}+1-\alpha_\sgn\le1-\alpha_*<0 ,
\]
because $1/\alpha_*-1/\alpha_\sgn=(\alpha_\sgn-\alpha_*)/(\alpha_*\alpha_\sgn)\le\alpha_\sgn-\alpha_*$,
and $B_NT_\sgn(N)\to0$ by Lemma~\ref{lem:rv}(v).
\end{proof}
